\section{Domain-Specific Figures and Evidence Anchors}

The following sections extend the comparison to implementation choices within access, serving, retrieval, and retention.

\subsection{Access-Scheduling Design Space}

Figure~\ref{fig:supp-access-designspace} relates the decision to use or move state to the observations and recovery rules that make the action possible.

\begin{figure*}[t]
\centering
\includegraphics[width=0.94\textwidth]{figures/pdf/access_scheduling_designspace.pdf}
\caption{Access-scheduling design space for locality, legality, transfer, and recovery.}
\label{fig:supp-access-designspace}
\end{figure*}

An authorization decision needs a visible version or progress boundary. A placement decision needs the cost of the next access. Recovery must preserve both the authority to act and the position from which work resumes.

Classical consistency and concurrency systems such as TicToc, Percolator, and RACE show that timestamping, metadata layout, and remote coordination are already control surfaces over state, even before modern serving runtimes made the same pattern explicit~\cite{yu2016tictoc,peng2010percolator,zuo2021race}. Systems such as CheckFreq, ChecknRun, PipeFill, and the newer runtime-monitoring lines show that the access problem is not just placement, but also when the runtime may safely observe, checkpoint, and resume without invalidating later actions~\cite{mohan2021checkfreq,eisenman2022checknrun,arfeen2024pipefill,zhou2024neomem}. Tooling for dynamic reconfiguration and transfer, including Caerus, MORPH-style recovery, and queue-aware stream control, matters because it keeps the state-ownership question visible instead of treating movement as a mere implementation detail~\cite{zhang2021caerus,mao2023morphstream,fu2024autocheck,gao2025deck}.

The same lineage also includes systems that turn placement pressure into explicit runtime policy rather than hidden implementation cost. Resource-multiplexing and scheduler-centric systems such as G10, resource multiplexing, PFS-style reconfiguration, and workload-sensitive monitoring show that queueing, hot-state exposure, and placement are inseparable once the runtime manages shared state across stages or tenants~\cite{jiang2024megascale,he2025resourcemultiplexing,domingo2021pfsck,jiang2025shadowviews}. Transactional dataflow and distributed-storage systems such as Flash-style state flow, Cloud-style buffering, and multi-stage transfer control reinforce the same point: the access problem is always also a visibility problem and a handoff problem~\cite{fu2022sfs,fu2024autocheck,gao2025deck,li2020pegasus}. That is why the survey keeps access and scheduling as a first-class axis rather than folding it into generic runtime management.

Earlier stream systems exposed where events waited and let runtimes change operator assignments~\cite{zhang2016corun,zhang2016elastic,zhang2019briskstream}. Checkpoint frequency and restoration add a time dimension: the state observed before a failure may differ from the last state that can be restored~\cite{mohan2021checkfreq,eisenman2022checknrun}.

Shared-device controllers combine residency, contention, and placement~\cite{cai2025moc,feng2025optimus,fsmoe25,wang2025wlbllm,huang2025tigon,zheng2024vmcu,tabatabai2024fbmm}. Storage disaggregation makes a similar tradeoff between flexible placement and access to state across a link~\cite{vuppalapati2020snowflake,zhou2024neomem}.

DRTMH-style transaction scheduling and the state-machine tradition remind us that legality over mutable state was a control problem long before modern runtimes named it that way~\cite{ren2015drtmh,schneider1990state}. EndMyth-style stream repair work shows the same issue from the failure side: access, recovery, and trust in observations are inseparable once the system can replay or resume~\cite{zamanian2017endmyth}.

\subsection{Serving Lifecycle Control}

When a request arrives, a serving system must decide whether an existing KV prefix can be reused and where to retrieve it.

SplitZip, CacheFlow, and ThunderServe make transfer cost and cloud heterogeneity explicit; ChandGPT-style reuse systems, Jenga, Prism, and Aegaeon make multi-model or multi-tenant residency explicit; Spinfer, FlexGen, and related offload designs make memory pressure and phase placement explicit; HeterInfer, MoE-Lightning, and resource multiplexing make the hardware/topology asymmetry explicit~\cite{guo2026splitzip,nian2026cacheflow,jiang2025thunderserve,gao2025weaver,zhang2025jenga,yu2025prism,fan2025spinfer,sheng2023flexgen,chen2025heteroinfer,cao2025moelightning,he2025resourcemultiplexing}.

PagedAttention, Orca, and AlpaServe explain why batching and scheduling were already state-sensitive before the latest disaggregation wave; SlimPipe, What-If Stragglers, and Pumped trace-based systems explain why scheduling debt becomes visible only after phase or hardware asymmetry is modeled; output-length prediction and cache-aware routing systems show why admission and future-memory forecasting belong in the same discussion as restoration and transfer~\cite{kwon2023pagedattention,yu2022orca,li2023alpaserve,liu2025slimpipe,lin2025whatifstragglers,gong2025pastfuture,jiang2025shadowviews}.

CacheSlide and SpecInfer-style speculative reuse make legality and verification a precondition for cross-position or tree-structured reuse; CachedAttention and HCache expose restoration cost across conversation turns; and MELL, Libra, fairness-aware serving, LayerKV, and Prefill-Only show that GPU residency, request partitioning, fairness, layer-wise allocation, and prefill admission determine which cached state remains reusable later~\cite{cacheslide2025,miao2024specinfer,gao2025cachedattention,gao2025hcache,liu2025mell,ruan2026libra,sheng2024fairnessserving,xiong2024layerkv,du2025prefillonly}. These systems make reuse, restoration, contention, and initialization part of the same serving control problem rather than separate local optimizations.

Application-centric serving such as Parrot shows that workflow-level state and prompt DAGs can dominate placement decisions even before the cache story starts~\cite{lin2024parrot}. Tree-shaped speculative and group-wise systems such as FastTree and InstAttention show that reuse and scheduling depend on structure in the request stream, not just on raw throughput targets~\cite{pan2025fasttree,pan2025instattention}. Decentralized and overlay-based systems such as PlanetServe and GLLM-adjacent coordination lines show that cache reuse, routing, and trust become intertwined once the scheduler is no longer centralized~\cite{fang2026planetserve}. The same cluster also includes systems such as Infinity-style eviction, Duplex-style cross-turn reuse, and world-load balancing variants, which are best understood as refinements of lifecycle control rather than new top-level subfields~\cite{infinigen2024,yun2024duplex,wang2025wlbllm}.

Prefill-optimized training-serving crossover systems such as MEPipe show that the serving boundary is often governed by how pipeline slack is distributed, not just by the cache itself~\cite{sun2025mepipe}. Copy-aware or clone-style state movement lines show the same thing from a different angle: once request state is shareable, the runtime has to decide whether cloning, relocating, or reusing is cheaper than recomputing~\cite{tian2025clone}.

Keeping an adapter resident saves a later load~\cite{xiang2025aegaeon,wu2024dlora}. Reusing a prompt prefix saves computation only when the stored state is valid for the new request~\cite{gim2024promptcache}.


\begin{figure*}[t]
\centering
\includegraphics[width=0.94\textwidth]{figures/pdf/kvcache_lifecycle.pdf}
\caption{KV-cache lifecycle.}
\label{fig:supp-kvcache-lifecycle}
\end{figure*}

The lifecycle figure makes the control seams concrete. Admit decides whether a request is allowed to create new short-lived state at all. Place decides where that state lives and therefore which queues or buses it will burden next. Mutate covers decode-time growth, prefix edits, and other in-place evolution of live serving state. Expose/transfer covers whether prefix, adapter, or session-local state may be reused safely and when it may move across workers or phases. Compact covers fragmentation control and structural cleanup. Evict/reclaim covers what can be released without hiding future restoration debt and how that debt is later charged back to the original SLO.

The following evidence is especially useful in this subsection:
\begin{itemize}[leftmargin=*]
\item predictive-serving families grouped by the debt they forecast;
\item legality conditions for prefix, branch, and multi-call reuse;
\item temporal-lifecycle policies over expose/transfer, evict/reclaim, warm-start, and fairness debt;
\item clearer distinctions between allocator realism, transfer scheduling, and ownership semantics.
\end{itemize}

\subsection{Retrieval and Retention Governance}

Retrieval and continual learning both depend on state that outlives the current operation. Indexes expose versions and pending mutations that affect what later queries can retrieve~\cite{singh2021freshdiskann,wang2021milvus,guo2022manu,asai2024selfrag,zhang2026flowrag,wang2026candor}. Continual learners must also account for the state they keep, including protected parameters, retained examples, compressed features, and the available memory budget~\cite{kirkpatrick2017ewc,rebuffi2017icarl,aljundi2019mir,prabhu2020gdumb,hayes2020remind,zhou2025ferret,li2026streamfp}.

Retrieval systems expose at least four trigger layers: request-time invocation, index-time maintenance, publication-time exposure, and placement-time residency. Self-RAG changes whether external memory should be consulted at all; FlowRAG changes when retriever state should be refreshed before reuse quality decays too far; CANDOR-Bench and FreshDiskANN expose when update churn, deletion debt, or local repair have already pushed the index away from its intended operating region; Milvus and Manu expose when partially refreshed state may be published safely under concurrent updates; Rummy and composable-memory search expose when movement or placement is justified by downstream quality gain instead of raw lookup throughput~\cite{asai2024selfrag,zhang2026flowrag,wang2026candor,singh2021freshdiskann,wang2021milvus,guo2022manu,zhang2024rummy,quinn2025accelerating}. The comparison shows that retrieval quality is a lifecycle property, not merely an encoder property.

Retention systems expose an analogous trigger hierarchy over one bounded store. EWC protects already-committed parameter state; iCaRL, GDumb, and StreamFP govern what enters the scarce retained set; GEM, A-GEM, MIR, and DER govern when stored state should re-enter the update path; REMIND governs whether retained signal is compressed and later reconstructed; Ferret governs what happens when the budget itself contracts under live update pressure~\cite{kirkpatrick2017ewc,rebuffi2017icarl,prabhu2020gdumb,li2026streamfp,lopez2017gem,chaudhry2019agem,aljundi2019mir,buzzega2020der,hayes2020remind,zhou2025ferret}. Read this way, the continual-learning literature matters to the survey not because it supplies another application area, but because it provides one of the clearest examples of how protect, admit, replay, compress, and budget-shrink decisions compose over time.

For retrieval, a service policy must say whether a stale local copy is permitted and when a changed version becomes queryable. For retention, the policy must say which evidence can be compressed or removed under a tighter budget. Subsequent quality measurements reveal the cost of that decision.

Dynamic ANN and index-maintenance lines such as HNSW, DiskANN, Milvus, and graph- or shard-centric systems make clear that freshness, shard publication, and rebuild debt are all part of one operational budget~\cite{malkov2020hnsw,singh2021freshdiskann,wang2021milvus,guo2022manu}. Retrieval-augmented and memory-augmented systems such as DPR, ColBERT, REALM, RAG, ATLAS, HippoRAG, and Raptor show the same phenomenon at the model boundary: the runtime must manage what gets exposed, refreshed, abstracted, and reused rather than treating memory as a static auxiliary index~\cite{karpukhin2020dpr,khattab2020colbert,guu2020realm,lewis2020rag,izacard2022atlas,wang2024hipporag,sarthi2024raptor}.



\begin{figure*}[t]
\centering
\includegraphics[width=0.94\textwidth]{figures/pdf/retention_strategy_taxonomy.pdf}
\caption{Retention-governance taxonomy.}
\label{fig:supp-retention-taxonomy}
\end{figure*}

Figure~\ref{fig:supp-retention-taxonomy} distinguishes decisions by the work they affect. Admission changes future coverage. Replay determines when retained state re-enters training. Compression changes what replay can recover, and budget contraction determines which items or guarantees are lost first.

\subsection{Long-Horizon Evolution Anchors}

SuperNeurons and related memory-reuse systems show that activation liveness, recomputation, and offload policy already form a managed state loop in deep-learning training, not just in inference-side serving~\cite{wang2018superneurons}. Bubble-filling and recomputation-aware training systems such as Obscura show the same thing from another angle: when state is expensive to keep live, the runtime must decide which debt to shift into pipeline slack and which debt to pay immediately~\cite{huang2025obscura}. That is the same governance pattern that later appears in retention stores, but with a different temporal horizon and a different type of reuse value.

\subsection{Blueprint and Fault-Model Boundary Conditions}

When mechanisms share state, each needs to declare the object it acts on, what it can observe, when its decision takes precedence, and what it may change. The recovery rules must preserve the invariants that other mechanisms rely on.

Recovery needs to restore state with the metadata used to validate it in steady state. A restored cache page needs its compatibility information. A recovered index segment needs the version that queries may observe.

Crash recovery with replay, a network partition, and delayed telemetry leave different state available to the controller~\cite{chandy1985snapshot,corbett2012spanner,ongaro2014raft,mao2023morphstream,zhao2024recovery,wang2026candor,zhang2026flowrag}. State the failure model before reasoning about whether a replica can publish a version, a new worker can acquire write authority, or a retention tier can discard an object.

A segment's visibility rules should name which queries may use it after a rebuild, deletion repair, or compaction. A retained object's rules should identify which fidelity, provenance, or availability conditions still hold after it is compressed or moved.

\subsection{Foundations Figures and Boundary Material}

A study enters the core comparison when it exposes a runtime decision over an identifiable state object and connects the decision to a service outcome. The overview figures map those decisions and their connection to the examples.

In state-machine replication, replicas apply an agreed sequence of commands to the same starting state~\cite{schneider1990state}. This makes the order of updates part of the correctness condition.

\inlinefigurewithcaption
  {width=0.9\linewidth}
  {figures/pdf/poster_core_overview.pdf}
  {Poster-style survey overview. It serves as a compact map of the survey's scope and argumentative flow.}
  {fig:supp-poster-overview}

Figure~\ref{fig:supp-running-examples} follows one change in workload in each service: account activity shifts, a long prompt arrives, or documents awaiting deletion accumulate. Read across each row from the state under pressure to the operation whose effect must be measured.

\inlinefigurewithcaption
  {width=0.96\linewidth}
  {figures/pdf/running_examples_triptych.pdf}
  {Running examples that recur across access, execution, and evolution.}
  {fig:supp-running-examples}
